\documentclass[11pt]{article}
\usepackage{amsmath, amssymb, amsthm}
\usepackage{geometry}
\usepackage{enumitem}
\usepackage{hyperref}

\geometry{a4paper, margin=1in}

\title{Comprehensive Logic and Proof Practice\\
\large 25 Multi-part Questions Derived from Core Concepts with Solutions}
\author{}
\date{}

\begin{document}

\maketitle

\tableofcontents

\newpage

\section*{Contents of Tested Concepts}
\addcontentsline{toc}{section}{Contents of Tested Concepts}
This custom contents page maps the core logical and set-theoretic concepts to the specific questions where they are tested across the 5 variants.

\begin{itemize}
    \item \textbf{Set Theory Fundamentals (Power Sets, Cardinality, Elements)}\dotfill V1-V5: Q1(i), Q1(ii)
    \item \textbf{Set Operations \& Cartesian Products}\dotfill V1-V5: Q1(iii)
    \item \textbf{Logic \& Truth Tables}\dotfill V1-V5: Q1(iv), Q3
    \item \textbf{Logical Equivalence, Simplification \& Negation}\dotfill V1-V5: Q1(v)
    \item \textbf{Quantified Statements (Evaluating \& Negating)}\dotfill V1-V5: Q1(vi), Q1(vii)
    \item \textbf{Vacuous \& Trivial Truths}\dotfill V1-V5: Q1(viii)
    \item \textbf{Set-Builder Notation \& Inequalities}\dotfill V1-V5: Q1(ix)
    \item \textbf{Infinite Indexed Collections \& Limits of Sets}\dotfill V1-V5: Q1(x)
    \item \textbf{Set Equality Proofs (Element Method)}\dotfill V1-V5: Q2
    \item \textbf{Direct Proofs \& Biconditionals}\dotfill V1-V3: Q4, Q5; V4: Q4; V5: Q4
    \item \textbf{Proof by Contrapositive}\dotfill V3: Q5, V5: Q5
    \item \textbf{Proof by Cases}\dotfill V4: Q5
\end{itemize}

\vspace{1cm}
\hrulefill

\newpage
\section{Variant 1: Fundamental Level (Easy)}

\subsection*{Question 1: Foundations of Sets and Logic}
\begin{enumerate}[label=(\roman*)]
    \item Let $A = \{a, b\}$. List the elements in $\mathcal{P}(A)$.\\
    \textbf{Solution:} The power set contains all subsets of $A$. \\
    $\mathcal{P}(A) = \{\emptyset, \{a\}, \{b\}, \{a, b\}\}$.

    \item Let $|A| = k$ and $|\mathcal{P}(A)| = 32$. Find $k$.\\
    \textbf{Solution:} The cardinality of a power set is $2^{|A|}$. \\
    $2^k = 32 \Rightarrow k = 5$.

    \item Let $A = \{1, 2\}$ and $B = \{3\}$. Find $|A \times B|$.\\
    \textbf{Solution:} For finite sets, $|A \times B| = |A| \cdot |B|$. \\
    $|A \times B| = 2 \cdot 1 = 2$.

    \item Complete the truth table for the expression $p \wedge q \Rightarrow p$.\\
    \textbf{Solution:} 
    \begin{center}
    \begin{tabular}{|c|c|c|c|}
    \hline
    $p$ & $q$ & $p \wedge q$ & $p \wedge q \Rightarrow p$ \\
    \hline
    T & T & T & T \\
    T & F & F & T \\
    F & T & F & T \\
    F & F & F & T \\
    \hline
    \end{tabular}
    \end{center}

    \item Simplify and eliminate the implication arrow: $\sim(p \Rightarrow q) \equiv$ \rule{3cm}{0.15mm} \\
    \textbf{Solution:} By the definition of implication, $p \Rightarrow q \equiv \sim p \vee q$. By De Morgan's Law, negating this gives $p \wedge \sim q$.

    \item Negate the following quantified statement: $\forall x \in \mathbb{R}, x > 0 \Rightarrow x^2 > 0$.\\
    \textbf{Solution:} $\exists x \in \mathbb{R}, x > 0 \wedge x^2 \leq 0$.

    \item Determine the truth value: $\exists x \in \mathbb{N}, x^2 - 4 = 0$.\\
    \textbf{Solution:} \textbf{True}. $x^2 - 4 = 0 \Rightarrow x = 2$ or $x = -2$. Since $2 \in \mathbb{N}$, a valid witness exists.

    \item Evaluate the truth value (T/F): Let $x \in \mathbb{R}$. If $x^2 < 0 \Rightarrow x = 5$.\\
    \textbf{Solution:} \textbf{True}. Since the square of any real number is non-negative, the hypothesis $x^2 < 0$ is always False. Thus, the implication is vacuously True.

    \item Let $A = \{x \in \mathbb{Z} : -2 < x \leq 1\}$. List the elements of $A$.\\
    \textbf{Solution:} The integers satisfying the inequality are $-1, 0, 1$. \\
    $A = \{-1, 0, 1\}$.

    \item Evaluate the infinite union of indexed sets: $\bigcup_{n=1}^\infty \left[0, 1-\frac{1}{n}\right]$.\\
    \textbf{Solution:} For $n=1$, the interval is $[0, 0] = \{0\}$. As $n \to \infty$, $1 - \frac{1}{n}$ approaches $1$ but never quite reaches it for any finite $n$. Thus, the union contains all real numbers from $0$ up to, but not including, $1$. \\
    Result: $[0, 1)$.
\end{enumerate}

\subsection*{Question 2: Set Equality Proof}
Let $A, B$, and $C$ be subsets of a universal set $U$. Prove using the element method that $A \cup (B \cap C) = (A \cup B) \cap (A \cup C)$.\\
\textbf{Solution:} 
\begin{proof}
Let $x$ be an arbitrary element in $U$. \\
$x \in A \cup (B \cap C)$ \\
$\Leftrightarrow x \in A \vee x \in (B \cap C)$ (Definition of Union) \\
$\Leftrightarrow x \in A \vee (x \in B \wedge x \in C)$ (Definition of Intersection) \\
$\Leftrightarrow (x \in A \vee x \in B) \wedge (x \in A \vee x \in C)$ (Distributive Law of Logic) \\
$\Leftrightarrow (x \in A \cup B) \wedge (x \in A \cup C)$ (Definition of Union) \\
$\Leftrightarrow x \in (A \cup B) \cap (A \cup C)$ (Definition of Intersection). \\
Since the biconditional holds for all $x$, the sets are equal.
\end{proof}

\subsection*{Question 3: Truth Tables and Tautologies}
Let $p, q$ be statements. Use truth tables to prove that $p \Rightarrow (p \vee q)$ is a tautology.\\
\textbf{Solution:}
\begin{center}
\begin{tabular}{|c|c|c|c|}
\hline
$p$ & $q$ & $p \vee q$ & $p \Rightarrow (p \vee q)$ \\
\hline
T & T & T & T \\
T & F & T & T \\
F & T & T & T \\
F & F & F & T \\
\hline
\end{tabular}
\end{center}
Since the final column consists entirely of True (T) values, the statement is a tautology.

\subsection*{Question 4: Biconditional Parity Proof}
Let $x \in \mathbb{Z}$. Prove: $x$ is odd $\Leftrightarrow x+1$ is even.\\
\textbf{Solution:}
\begin{proof}
We prove both directions. \\
($\Rightarrow$): Assume $x$ is odd. By definition, $\exists k \in \mathbb{Z}$ such that $x = 2k + 1$. Then $x+1 = (2k+1)+1 = 2k+2 = 2(k+1)$. Since $k+1 \in \mathbb{Z}$, $x+1$ is even. \\
($\Leftarrow$): Assume $x+1$ is even. By definition, $\exists m \in \mathbb{Z}$ such that $x+1 = 2m$. Then $x = 2m - 1 = 2m - 2 + 1 = 2(m-1) + 1$. Since $m-1 \in \mathbb{Z}$, $x$ is odd.
\end{proof}

\subsection*{Question 5: Direct Proof}
Let $x \in \mathbb{Z}$. Prove that if $x$ is odd, then $3x+2$ is odd.\\
\textbf{Solution:}
\begin{proof}
Assume $x$ is an odd integer. Then $x = 2k + 1$ for some integer $k$. We evaluate $3x + 2$: \\
$3x + 2 = 3(2k + 1) + 2 = 6k + 3 + 2 = 6k + 5$. \\
We rewrite this to show it has the form of an odd integer: \\
$6k + 5 = 6k + 4 + 1 = 2(3k + 2) + 1$. \\
Since $k$ is an integer, $3k + 2$ is an integer. Thus, $3x + 2$ is odd.
\end{proof}

\newpage
\section{Variant 2: Moderate Level (Standard Exam Difficulty)}

\subsection*{Question 1: Foundations of Sets and Logic}
\begin{enumerate}[label=(\roman*)]
    \item Let $A = \{\emptyset, 1\}$. List the elements in $\mathcal{P}(A)$.\\
    \textbf{Solution:} $\mathcal{P}(A) = \{\emptyset, \{\emptyset\}, \{1\}, \{\emptyset, 1\}\}$.

    \item Let $|A| = k$ and $|\mathcal{P}(A)|^2 - 8|\mathcal{P}(A)| = 0$. Find $k$.\\
    \textbf{Solution:} Let $y = |\mathcal{P}(A)| = 2^k$. The equation is $y^2 - 8y = 0 \Rightarrow y(y-8) = 0$. Since cardinality cannot be 0, $y = 8$. Thus, $2^k = 8 \Rightarrow k = 3$.

    \item Let $A = \{1\}, B = \{1, 2\}, C = \{1, 3\}$. Find $|(A \times B) \cup (A \times C)|$.\\
    \textbf{Solution:} By distributivity, $(A \times B) \cup (A \times C) = A \times (B \cup C)$. \\
    $B \cup C = \{1, 2, 3\}$, which has size 3. $|A| = 1$. The cardinality is $1 \cdot 3 = 3$.

    \item Complete the truth table for the expression $(p \Rightarrow q) \wedge \sim q$.\\
    \textbf{Solution:}
    \begin{center}
    \begin{tabular}{|c|c|c|c|c|}
    \hline
    $p$ & $q$ & $p \Rightarrow q$ & $\sim q$ & $(p \Rightarrow q) \wedge \sim q$ \\
    \hline
    T & T & T & F & F \\
    T & F & F & T & F \\
    F & T & T & F & F \\
    F & F & T & T & T \\
    \hline
    \end{tabular}
    \end{center}

    \item Simplify and eliminate the implication arrow: $\sim[(p \vee q) \Rightarrow r] \equiv$ \rule{3cm}{0.15mm}\\
    \textbf{Solution:} $\sim(A \Rightarrow B) \equiv A \wedge \sim B$. Let $A = p \vee q$ and $B = r$. The result is $(p \vee q) \wedge \sim r$.

    \item Negate the following quantified statement: $\exists x \in S, \forall y \in T, x+y = 0 \Rightarrow x = y$.\\
    \textbf{Solution:} $\forall x \in S, \exists y \in T, x+y = 0 \wedge x \neq y$.

    \item Determine the truth value: $\forall x \in \mathbb{R}, \exists y \in \mathbb{R}, x+y = 0$.\\
    \textbf{Solution:} \textbf{True}. For any real number $x$, we can select $y = -x$, which is also a real number, and $x + (-x) = 0$.

    \item Evaluate the truth value (T/F): If $0 = 1 \Rightarrow$ all integers are even.\\
    \textbf{Solution:} \textbf{True}. The hypothesis $0=1$ is False. A false hypothesis makes an implication vacuously True.

    \item Let $A = \{x \in \mathbb{Z} : |2x - 1| < 4\}$. List the elements of $A$.\\
    \textbf{Solution:} $|2x - 1| < 4 \Rightarrow -4 < 2x - 1 < 4 \Rightarrow -3 < 2x < 5 \Rightarrow -1.5 < x < 2.5$. The integers in this range are $-1, 0, 1, 2$. \\
    $A = \{-1, 0, 1, 2\}$.

    \item Evaluate the infinite intersection of indexed sets: $\bigcap_{n=1}^\infty \left(-\frac{1}{n}, \frac{1}{n}\right)$.\\
    \textbf{Solution:} As $n \to \infty$, the bounds shrink towards 0. The only number strictly greater than $-1/n$ and strictly less than $1/n$ for \textit{all} $n \geq 1$ is 0. \\
    Result: $\{0\}$.
\end{enumerate}

\subsection*{Question 2: Set Equality Proof}
Let $A$ and $B$ be subsets of a universal set $U$. Prove using the element method that $(A \cap B)^c = A^c \cup B^c$ (De Morgan's Law for Sets).\\
\textbf{Solution:}
\begin{proof}
Let $x \in U$. \\
$x \in (A \cap B)^c$ \\
$\Leftrightarrow \sim(x \in A \cap B)$ (Definition of Complement) \\
$\Leftrightarrow \sim(x \in A \wedge x \in B)$ (Definition of Intersection) \\
$\Leftrightarrow \sim(x \in A) \vee \sim(x \in B)$ (De Morgan's Law for Logic) \\
$\Leftrightarrow x \in A^c \vee x \in B^c$ (Definition of Complement) \\
$\Leftrightarrow x \in A^c \cup B^c$ (Definition of Union). \\
Since the biconditional is established, $(A \cap B)^c = A^c \cup B^c$.
\end{proof}

\subsection*{Question 3: Truth Tables and Tautologies}
Let $p, q$ be statements. Use truth tables to prove that $(p \Rightarrow q) \Leftrightarrow (\sim q \Rightarrow \sim p)$ is a tautology.\\
\textbf{Solution:}
\begin{center}
\begin{tabular}{|c|c|c|c|c|c|c|}
\hline
$p$ & $q$ & $\sim p$ & $\sim q$ & $p \Rightarrow q$ & $\sim q \Rightarrow \sim p$ & $(p \Rightarrow q) \Leftrightarrow (\sim q \Rightarrow \sim p)$ \\
\hline
T & T & F & F & T & T & T \\
T & F & F & T & F & F & T \\
F & T & T & F & T & T & T \\
F & F & T & T & T & T & T \\
\hline
\end{tabular}
\end{center}
The final column is all True, meaning the biconditional is a tautology. This verifies the contrapositive equivalence.

\subsection*{Question 4: Biconditional Parity Proof}
Let $x, y \in \mathbb{Z}$. Prove: $x$ and $y$ have the same parity $\Leftrightarrow x+y$ is even.\\
\textbf{Solution:}
\begin{proof}
We prove both directions. \\
($\Rightarrow$): Assume $x, y$ have the same parity. We have two cases:
\begin{itemize}
    \item Case 1: Both are even. $x = 2k, y = 2m$. $x+y = 2k+2m = 2(k+m)$. Even.
    \item Case 2: Both are odd. $x = 2k+1, y = 2m+1$. $x+y = 2k+1+2m+1 = 2k+2m+2 = 2(k+m+1)$. Even.
\end{itemize}
In both cases, $x+y$ is even. \\
($\Leftarrow$): We proceed by contrapositive. Assume $x$ and $y$ do NOT have the same parity (i.e., one is even, one is odd). Without loss of generality, let $x = 2k$ and $y = 2m+1$. \\
$x+y = 2k + 2m + 1 = 2(k+m) + 1$. \\
This is an odd integer, meaning $x+y$ is not even. This proves the converse.
\end{proof}

\subsection*{Question 5: Direct Proof}
Let $x \in \mathbb{Z}$. Prove that if $x$ is even, then $5x-3$ is odd.\\
\textbf{Solution:}
\begin{proof}
Assume $x$ is even. Then $x = 2k$ for some $k \in \mathbb{Z}$. \\
Substitute $2k$ into $5x - 3$: \\
$5(2k) - 3 = 10k - 3$. \\
To show this is odd, we manipulate the expression to the form $2m+1$: \\
$10k - 3 = 10k - 4 + 1 = 2(5k - 2) + 1$. \\
Let $m = 5k - 2$. Since $k \in \mathbb{Z}$, $m \in \mathbb{Z}$. Since $5x - 3 = 2m + 1$, it is an odd integer.
\end{proof}

\newpage
\section{Variant 3: Intermediate Level (Advanced Manipulations)}

\subsection*{Question 1: Foundations of Sets and Logic}
\begin{enumerate}[label=(\roman*)]
    \item Let $A = \{1, \{1, 2\}\}$. List the elements in $\mathcal{P}(A)$.\\
    \textbf{Solution:} $\mathcal{P}(A) = \{\emptyset, \{1\}, \{\{1, 2\}\}, \{1, \{1, 2\}\}\}$.

    \item Let $|A| = k$ and $2^{|A|} \cdot |\mathcal{P}(A)| = 64$. Find $k$.\\
    \textbf{Solution:} $2^k \cdot 2^k = 64 \Rightarrow 2^{2k} = 2^6 \Rightarrow 2k = 6 \Rightarrow k = 3$.

    \item Let $A = \{x, y\}$ and $B = \{y, z\}$. Find $|(A \times A) \cap (B \times B)|$.\\
    \textbf{Solution:} Using the Cartesian intersection identity: $(A \times A) \cap (B \times B) = (A \cap B) \times (A \cap B)$. \\
    $A \cap B = \{y\}$. Thus, $|\{y\} \times \{y\}| = 1 \cdot 1 = 1$.

    \item Complete the truth table for the expression $p \Rightarrow (q \Rightarrow p)$.\\
    \textbf{Solution:}
    \begin{center}
    \begin{tabular}{|c|c|c|c|}
    \hline
    $p$ & $q$ & $q \Rightarrow p$ & $p \Rightarrow (q \Rightarrow p)$ \\
    \hline
    T & T & T & T \\
    T & F & T & T \\
    F & T & F & T \\
    F & F & T & T \\
    \hline
    \end{tabular}
    \end{center}

    \item Simplify and eliminate the implication arrow: $\sim[\sim p \Rightarrow (q \wedge \sim r)] \equiv$ \rule{3cm}{0.15mm}\\
    \textbf{Solution:} $\sim(A \Rightarrow B) \equiv A \wedge \sim B$. \\
    Here, $A = \sim p$ and $B = q \wedge \sim r$. \\
    Result: $\sim p \wedge \sim(q \wedge \sim r) \equiv \sim p \wedge (\sim q \vee r)$.

    \item Negate the following statement: $\forall x \in \mathbb{R}, \exists y \in \mathbb{R}, (x>0 \wedge y<0) \Rightarrow xy < 0$.\\
    \textbf{Solution:} $\exists x \in \mathbb{R}, \forall y \in \mathbb{R}, x>0 \wedge y<0 \wedge xy \geq 0$.

    \item Determine the truth value: $\exists x \in \mathbb{Q}, x^2 = 2$.\\
    \textbf{Solution:} \textbf{False}. The only real solutions to $x^2=2$ are $\pm\sqrt{2}$, which are irrational numbers ($\notin \mathbb{Q}$).

    \item Evaluate the truth value (T/F): Let $n \in \mathbb{N}$. If $n < 0 \Rightarrow n^2 < 0$.\\
    \textbf{Solution:} \textbf{True}. The domain is natural numbers ($\mathbb{N} = \{1, 2, ...\}$). $n < 0$ is always False. A false hypothesis creates a vacuously true implication.

    \item Let $A = \{x \in \mathbb{Z} : x^2 \leq 10 \text{ and } x \text{ is odd}\}$. List the elements of $A$.\\
    \textbf{Solution:} The integers whose squares are $\leq 10$ are $\{-3, -2, -1, 0, 1, 2, 3\}$. The odd integers in this set are $\{-3, -1, 1, 3\}$. \\
    $A = \{-3, -1, 1, 3\}$.

    \item Evaluate the infinite union of indexed sets: $\bigcup_{n=1}^\infty \left[\frac{1}{n}, 2-\frac{1}{n}\right]$.\\
    \textbf{Solution:} For $n=1$, the interval is $[1, 1] = \{1\}$. As $n$ increases, the lower bound approaches $0$ (but never reaches it) and the upper bound approaches $2$ (but never reaches it). \\
    Result: $(0, 2)$.
\end{enumerate}

\subsection*{Question 2: Set Equality Proof}
Let $A, B$, and $C$ be sets. Prove using the element method that $A - (B \cup C) = (A - B) \cap (A - C)$.\\
\textbf{Solution:}
\begin{proof}
Let $x$ be an arbitrary element. \\
$x \in A - (B \cup C)$ \\
$\Leftrightarrow x \in A \wedge x \notin (B \cup C)$ (Definition of Set Difference) \\
$\Leftrightarrow x \in A \wedge \sim(x \in B \vee x \in C)$ (Definition of Union) \\
$\Leftrightarrow x \in A \wedge (\sim(x \in B) \wedge \sim(x \in C))$ (De Morgan's Law for Logic) \\
$\Leftrightarrow (x \in A \wedge x \notin B) \wedge (x \in A \wedge x \notin C)$ (Idempotent and Distributive Laws) \\
$\Leftrightarrow x \in (A - B) \wedge x \in (A - C)$ (Definition of Set Difference) \\
$\Leftrightarrow x \in (A - B) \cap (A - C)$ (Definition of Intersection). \\
Thus, $A - (B \cup C) = (A - B) \cap (A - C)$.
\end{proof}

\subsection*{Question 3: Truth Tables and Tautologies}
Let $p, q$ be statements. Use truth tables to prove that $((p \Rightarrow q) \wedge p) \Rightarrow q$ is a tautology. (Modus Ponens)\\
\textbf{Solution:}
\begin{center}
\begin{tabular}{|c|c|c|c|c|}
\hline
$p$ & $q$ & $p \Rightarrow q$ & $(p \Rightarrow q) \wedge p$ & $((p \Rightarrow q) \wedge p) \Rightarrow q$ \\
\hline
T & T & T & T & T \\
T & F & F & F & T \\
F & T & T & F & T \\
F & F & T & F & T \\
\hline
\end{tabular}
\end{center}
The final column is all True; thus, Modus Ponens is a tautology.

\subsection*{Question 4: Biconditional Parity Proof}
Let $x \in \mathbb{Z}$. Prove: $x^2$ is even $\Leftrightarrow x$ is even.\\
\textbf{Solution:}
\begin{proof}
We prove both directions. \\
($\Leftarrow$): Assume $x$ is even. $x = 2k$. Then $x^2 = (2k)^2 = 4k^2 = 2(2k^2)$. Since $2k^2 \in \mathbb{Z}$, $x^2$ is even. \\
($\Rightarrow$): We prove this direction by contrapositive. The contrapositive statement is: If $x$ is odd, then $x^2$ is odd. \\
Assume $x$ is odd. Then $x = 2m + 1$ for some $m \in \mathbb{Z}$. \\
$x^2 = (2m + 1)^2 = 4m^2 + 4m + 1 = 2(2m^2 + 2m) + 1$. \\
Since $2m^2+2m \in \mathbb{Z}$, $x^2$ is odd. Thus, the contrapositive (and the original forward direction) is proven.
\end{proof}

\subsection*{Question 5: Contrapositive Proof}
Let $x \in \mathbb{Z}$. Prove that if $7x-8$ is even, then $x$ is even.\\
\textbf{Solution:}
\begin{proof}
We proceed by contrapositive. Assume $x$ is odd. We must show $7x - 8$ is odd. \\
Since $x$ is odd, $x = 2k + 1$ for $k \in \mathbb{Z}$. \\
$7x - 8 = 7(2k + 1) - 8 = 14k + 7 - 8 = 14k - 1$. \\
To match the form of an odd integer, we rewrite: \\
$14k - 1 = 14k - 2 + 1 = 2(7k - 1) + 1$. \\
Since $7k - 1 \in \mathbb{Z}$, $7x - 8$ is odd. The contrapositive is proven.
\end{proof}

\newpage
\section{Variant 4: Advanced Level (Edge Cases and Multi-Step Logic)}

\subsection*{Question 1: Foundations of Sets and Logic}
\begin{enumerate}[label=(\roman*)]
    \item Let $A = \{\emptyset, \{\emptyset\}\}$. List the elements in $\mathcal{P}(A)$.\\
    \textbf{Solution:} $\mathcal{P}(A) = \{\emptyset, \{\emptyset\}, \{\{\emptyset\}\}, \{\emptyset, \{\emptyset\}\}\}$.

    \item Let $|A| = 2$ and $|B| = m$. If $|\mathcal{P}(A \times B)| = 256$, find $m$.\\
    \textbf{Solution:} $|A \times B| = |A||B| = 2m$. The power set size is $2^{2m} = 256$. Since $2^8 = 256$, $2m = 8 \Rightarrow m = 4$.

    \item Let $A = \{1, 2, 3\}$ and $B = \{2, 3, 4\}$. Find $|(A \times B) - (B \times A)|$.\\
    \textbf{Solution:} $|A \times B| = 3 \cdot 3 = 9$. The overlapping elements are $(A \times B) \cap (B \times A) = (A \cap B) \times (A \cap B)$. $A \cap B = \{2, 3\}$, so there are $2 \cdot 2 = 4$ shared ordered pairs. \\
    Thus, $9 - 4 = 5$.

    \item Complete the truth table for the expression $(p \vee q) \wedge \sim p \Rightarrow q$.\\
    \textbf{Solution:}
    \begin{center}
    \begin{tabular}{|c|c|c|c|c|c|}
    \hline
    $p$ & $q$ & $p \vee q$ & $\sim p$ & $(p \vee q) \wedge \sim p$ & $((p \vee q) \wedge \sim p) \Rightarrow q$ \\
    \hline
    T & T & T & F & F & T \\
    T & F & T & F & F & T \\
    F & T & T & T & T & T \\
    F & F & F & T & F & T \\
    \hline
    \end{tabular}
    \end{center}

    \item Simplify and eliminate the implication arrow: $\sim[(p \Rightarrow q) \Rightarrow (p \wedge r)] \equiv$ \rule{3cm}{0.15mm}\\
    \textbf{Solution:} $\sim(A \Rightarrow B) \equiv A \wedge \sim B$. \\
    $(p \Rightarrow q) \wedge \sim(p \wedge r) \equiv (\sim p \vee q) \wedge (\sim p \vee \sim r)$.

    \item Negate the formal limit definition: $\forall \epsilon > 0, \exists \delta > 0, |x-a|<\delta \Rightarrow |f(x)-L|<\epsilon$.\\
    \textbf{Solution:} Domain definitions do not change sign during negation. \\
    $\exists \epsilon > 0, \forall \delta > 0, |x-a|<\delta \wedge |f(x)-L| \geq \epsilon$.

    \item Determine the truth value: $\forall x \in \mathbb{Z}, \exists y \in \mathbb{Z}, x-y \text{ is odd}$.\\
    \textbf{Solution:} \textbf{True}. If $x$ is even, pick $y=1$ (even $-$ odd $=$ odd). If $x$ is odd, pick $y=2$ (odd $-$ even $=$ odd).

    \item Evaluate the truth value (T/F): Let $x,y \in \mathbb{R}$. If $x^2+y^2 < 0 \Rightarrow x+y > 10$.\\
    \textbf{Solution:} \textbf{True}. $x^2+y^2 < 0$ is impossible for real numbers. A false hypothesis yields a vacuously true implication.

    \item Let $A = \{x \in \mathbb{Z} : x \equiv 1 \pmod 3 \text{ and } -5 \leq x \leq 5\}$. List the elements of $A$.\\
    \textbf{Solution:} The numbers equal to $3k+1$ within the bounds are $-5, -2, 1, 4$. \\
    $A = \{-5, -2, 1, 4\}$.

    \item Evaluate the infinite intersection of indexed sets: $\bigcap_{n=1}^\infty \left[0, \frac{1}{n}\right)$.\\
    \textbf{Solution:} As $n \to \infty$, $1/n$ approaches 0. Since the interval is open at $1/n$, 0 is the only element that satisfies $0 \leq x < 1/n$ for all integer values $n \geq 1$. \\
    Result: $\{0\}$.
\end{enumerate}

\subsection*{Question 2: Subset and Union Proof}
Let $A$ and $B$ be sets. Prove using the element method that if $A \subseteq B$, then $A \cup B = B$.\\
\textbf{Solution:}
\begin{proof}
We must prove $A \cup B \subseteq B$ and $B \subseteq A \cup B$. \\
1) Show $B \subseteq A \cup B$: Let $x \in B$. By the definition of union, $x \in A \vee x \in B$. Since $x \in B$ is true, $x \in A \cup B$. \\
2) Show $A \cup B \subseteq B$: Let $x \in A \cup B$. Then $x \in A$ or $x \in B$. 
\begin{itemize}
    \item If $x \in B$, we are done.
    \item If $x \in A$, since $A \subseteq B$ by our hypothesis, $x \in B$.
\end{itemize}
In all cases, $x \in B$. Thus $A \cup B \subseteq B$. \\
Since both inclusions hold, $A \cup B = B$.
\end{proof}

\subsection*{Question 3: Truth Tables and Logical Equivalence}
Let $p, q, r$ be statements. Use a truth table to show that $p \wedge (q \vee r) \Leftrightarrow (p \wedge q) \vee (p \wedge r)$.\\
\textbf{Solution:}
\begin{center}
\begin{tabular}{|c|c|c|c|c|c|c|c|}
\hline
$p$ & $q$ & $r$ & $q \vee r$ & $p \wedge (q \vee r)$ & $p \wedge q$ & $p \wedge r$ & $(p \wedge q) \vee (p \wedge r)$ \\
\hline
T & T & T & T & \textbf{T} & T & T & \textbf{T} \\
T & T & F & T & \textbf{T} & T & F & \textbf{T} \\
T & F & T & T & \textbf{T} & F & T & \textbf{T} \\
T & F & F & F & \textbf{F} & F & F & \textbf{F} \\
F & T & T & T & \textbf{F} & F & F & \textbf{F} \\
F & T & F & T & \textbf{F} & F & F & \textbf{F} \\
F & F & T & T & \textbf{F} & F & F & \textbf{F} \\
F & F & F & F & \textbf{F} & F & F & \textbf{F} \\
\hline
\end{tabular}
\end{center}
Columns 5 and 8 are identical, proving the logical equivalence.

\subsection*{Question 4: Biconditional Parity Proof}
Let $x, y \in \mathbb{Z}$. Prove: $xy$ is odd $\Leftrightarrow x$ is odd and $y$ is odd.\\
\textbf{Solution:}
\begin{proof}
($\Leftarrow$): Assume $x$ is odd and $y$ is odd. $x = 2k+1, y = 2m+1$. \\
$xy = (2k+1)(2m+1) = 4km + 2k + 2m + 1 = 2(2km + k + m) + 1$. \\
Since $2km+k+m \in \mathbb{Z}$, $xy$ is odd. \\
($\Rightarrow$): We prove by contrapositive. Assume it is not the case that both $x$ and $y$ are odd. By De Morgan's Law, $x$ is even OR $y$ is even. \\
Case 1: $x$ is even ($x=2k$). Then $xy = (2k)y = 2(ky)$. Since $ky \in \mathbb{Z}$, $xy$ is even. \\
Case 2: $y$ is even ($y=2m$). Then $xy = x(2m) = 2(xm)$. Since $xm \in \mathbb{Z}$, $xy$ is even. \\
In all cases of the negated conclusion, $xy$ is even (not odd). This proves the forward implication.
\end{proof}

\subsection*{Question 5: Proof by Cases}
Let $n \in \mathbb{Z}$. Prove by cases that the expression $n^2 - 3n + 9$ is always odd.\\
\textbf{Solution:}
\begin{proof}
We partition the integers into two cases based on parity. \\
\textbf{Case 1: $n$ is even.} Let $n = 2k$. \\
$n^2 - 3n + 9 = (2k)^2 - 3(2k) + 9 = 4k^2 - 6k + 8 + 1 = 2(2k^2 - 3k + 4) + 1$. \\
Since $2k^2-3k+4 \in \mathbb{Z}$, the expression is odd. \\
\textbf{Case 2: $n$ is odd.} Let $n = 2k+1$. \\
$n^2 - 3n + 9 = (2k+1)^2 - 3(2k+1) + 9 = 4k^2 + 4k + 1 - 6k - 3 + 9 = 4k^2 - 2k + 7$. \\
We rewrite this as: $4k^2 - 2k + 6 + 1 = 2(2k^2 - k + 3) + 1$. \\
Since $2k^2 - k + 3 \in \mathbb{Z}$, the expression is odd. \\
Because $n^2 - 3n + 9$ is odd in all exhaustive cases, the theorem is proven.
\end{proof}

\newpage
\section{Variant 5: Challenge Level (Abstract Synthesis)}

\subsection*{Question 1: Foundations of Sets and Logic}
\begin{enumerate}[label=(\roman*)]
    \item List the elements in $\mathcal{P}(\mathcal{P}(\{\emptyset\}))$.\\
    \textbf{Solution:} $\mathcal{P}(\{\emptyset\}) = \{\emptyset, \{\emptyset\}\}$. The power set of this 2-element set has 4 elements: \\
    $\{\emptyset, \{\emptyset\}, \{\{\emptyset\}\}, \{\emptyset, \{\emptyset\}\}\}$.

    \item Let $|A| = k$. Find $k$ if $(|\mathcal{P}(A)|)! = 24$ (where $!$ denotes the factorial function).\\
    \textbf{Solution:} $4! = 24$, so $|\mathcal{P}(A)| = 4$. Since $2^k = 4$, $k = 2$.

    \item Let $S = \{1, 2, 3, 4\}$ be partitioned by the set $P = \{\{1, 2\}, \{3\}, C\}$. Find the set $C$.\\
    \textbf{Solution:} Partitions are mutually exclusive and exhaustive. The missing element to complete the union of $S$ is 4. Thus, $C = \{4\}$.

    \item Complete the truth table for the expression $(p \Leftrightarrow q) \Rightarrow (p \wedge q)$.\\
    \textbf{Solution:}
    \begin{center}
    \begin{tabular}{|c|c|c|c|c|}
    \hline
    $p$ & $q$ & $p \Leftrightarrow q$ & $p \wedge q$ & $(p \Leftrightarrow q) \Rightarrow (p \wedge q)$ \\
    \hline
    T & T & T & T & T \\
    T & F & F & F & T \\
    F & T & F & F & T \\
    F & F & T & F & F \\
    \hline
    \end{tabular}
    \end{center}

    \item Simplify and eliminate the implication arrow: $\sim[\sim(p \Rightarrow \sim q) \vee (r \Rightarrow p)] \equiv$ \rule{3cm}{0.15mm}\\
    \textbf{Solution:} 
    $\sim[\sim(\sim p \vee \sim q) \vee (\sim r \vee p)]$ \\
    $\equiv \sim[(p \wedge q) \vee (\sim r \vee p)]$ \\
    By De Morgan's Law: $\sim(p \wedge q) \wedge \sim(\sim r \vee p)$ \\
    $\equiv (\sim p \vee \sim q) \wedge (r \wedge \sim p)$ \\
    By commutativity and associativity, factor out $\sim p$: \\
    $\equiv \sim p \wedge r \wedge (\sim p \vee \sim q) \equiv \sim p \wedge r$.

    \item Negate the uniqueness quantifier statement (conceptually): $\exists! x \in \mathbb{R}, x^2 = 0$.\\
    \textbf{Solution:} $\exists! x, P(x) \equiv \exists x (P(x) \wedge \forall y(P(y) \Rightarrow y=x))$. \\
    Negating this gives: $\forall x (\sim P(x) \vee \exists y(P(y) \wedge y \neq x))$. \\
    In words: Either no real number satisfies $x^2=0$, or there are at least two distinct real numbers that satisfy it.

    \item Determine the truth value: $\exists x \in \mathbb{R}, \forall y \in \mathbb{R}, xy = 0$.\\
    \textbf{Solution:} \textbf{True}. Let $x = 0$. For any $y \in \mathbb{R}$, $0 \cdot y = 0$.

    \item Evaluate the truth value (T/F): Let $S$ be the set of odd integers. $\forall x \in S$, if $x$ is even, then $x^2 = 4$.\\
    \textbf{Solution:} \textbf{True}. Because the domain $S$ only contains odd integers, the condition "$x$ is even" is always False. Vacuously True.

    \item Let $A = \{x \in \mathbb{Z} : x^2 < x\}$. List the elements of $A$.\\
    \textbf{Solution:} $x^2 - x < 0 \Rightarrow x(x-1) < 0 \Rightarrow 0 < x < 1$. There are no integers between 0 and 1. \\
    $A = \emptyset$.

    \item Evaluate the limit supremum of the sequence of sets: $\limsup_{n\to\infty} A_n$ where $A_n = [0, 2 - (-1)^n]$.\\
    \textbf{Solution:} Limit supremum represents elements that appear in infinitely many sets. \\
    For even $n$, $A_n = [0, 1]$. For odd $n$, $A_n = [0, 3]$. The elements in $[0, 3]$ appear infinitely many times (every odd step). \\
    Result: $[0, 3]$.
\end{enumerate}

\subsection*{Question 2: Set Equality Proof}
Let $A, B$, and $C$ be sets. Prove using the element method that $A \times (B \cap C) = (A \times B) \cap (A \times C)$.\\
\textbf{Solution:}
\begin{proof}
Let $(x,y)$ be an arbitrary ordered pair. \\
$(x,y) \in A \times (B \cap C)$ \\
$\Leftrightarrow x \in A \wedge y \in (B \cap C)$ (Definition of Cartesian Product) \\
$\Leftrightarrow x \in A \wedge (y \in B \wedge y \in C)$ (Definition of Intersection) \\
$\Leftrightarrow (x \in A \wedge y \in B) \wedge (x \in A \wedge y \in C)$ (Idempotent Law for $x \in A$) \\
$\Leftrightarrow (x,y) \in (A \times B) \wedge (x,y) \in (A \times C)$ (Definition of Cartesian Product) \\
$\Leftrightarrow (x,y) \in (A \times B) \cap (A \times C)$ (Definition of Intersection). \\
Thus, $A \times (B \cap C) = (A \times B) \cap (A \times C)$.
\end{proof}

\subsection*{Question 3: Truth Tables and Tautologies}
Let $p, q, r$ be statements. Use truth tables to prove the Law of Syllogism is a tautology: $((p \Rightarrow q) \wedge (q \Rightarrow r)) \Rightarrow (p \Rightarrow r)$.\\
\textbf{Solution:}
\begin{center}
\begin{tabular}{|c|c|c|c|c|c|c|c|}
\hline
$p$ & $q$ & $r$ & $p \Rightarrow q$ & $q \Rightarrow r$ & $(p \Rightarrow q) \wedge (q \Rightarrow r)$ & $p \Rightarrow r$ & Final \\
\hline
T & T & T & T & T & T & T & \textbf{T} \\
T & T & F & T & F & F & F & \textbf{T} \\
T & F & T & F & T & F & T & \textbf{T} \\
T & F & F & F & T & F & F & \textbf{T} \\
F & T & T & T & T & T & T & \textbf{T} \\
F & T & F & T & F & F & T & \textbf{T} \\
F & F & T & T & T & T & T & \textbf{T} \\
F & F & F & T & T & T & T & \textbf{T} \\
\hline
\end{tabular}
\end{center}
Since the Final column evaluates exclusively to True, it is a tautology.

\subsection*{Question 4: Biconditional Parity Proof}
Let $x,y \in \mathbb{Z}$. Prove: $3x + 5y$ is even $\Leftrightarrow x$ and $y$ are of the same parity.\\
\textbf{Solution:}
\begin{proof}
($\Leftarrow$): Assume $x,y$ share the same parity. \\
Case 1: Both even ($x=2k, y=2m$). $3(2k) + 5(2m) = 6k + 10m = 2(3k + 5m)$, which is even. \\
Case 2: Both odd ($x=2k+1, y=2m+1$). $3(2k+1) + 5(2m+1) = 6k + 3 + 10m + 5 = 6k + 10m + 8 = 2(3k + 5m + 4)$, which is even. \\
($\Rightarrow$): We prove by contrapositive. Assume $x,y$ are of opposite parity. \\
Case 1: $x$ even ($2k$), $y$ odd ($2m+1$). $3(2k) + 5(2m+1) = 6k + 10m + 5 = 2(3k+5m+2)+1$. Odd. \\
Case 2: $x$ odd ($2k+1$), $y$ even ($2m$). $3(2k+1) + 5(2m) = 6k + 3 + 10m = 2(3k+5m+1)+1$. Odd. \\
Since opposite parity always leads to $3x+5y$ being odd, the biconditional is proven.
\end{proof}

\subsection*{Question 5: Contrapositive Logic}
Let $x,y \in \mathbb{Z}$. Prove by contrapositive: If $xy$ and $x+y$ are even, then both $x$ and $y$ are even.\\
\textbf{Solution:}
\begin{proof}
The contrapositive statement is: "If it is not the case that both $x$ and $y$ are even, then $xy$ is odd OR $x+y$ is odd." \\
Assume it is not the case that both $x$ and $y$ are even. By De Morgan's Laws, this means $x$ is odd OR $y$ is odd. We proceed by cases:
\begin{itemize}
    \item \textbf{Case 1: Both $x$ and $y$ are odd.} Let $x=2k+1$ and $y=2m+1$. \\
    Then $xy = (2k+1)(2m+1) = 4km+2k+2m+1 = 2(2km+k+m)+1$, which is odd. This satisfies the "OR" condition of the conclusion.
    \item \textbf{Case 2: Exactly one is odd, the other is even.} Without loss of generality, let $x=2k+1$ and $y=2m$. \\
    Then $x+y = (2k+1) + 2m = 2k+2m+1 = 2(k+m)+1$, which is odd. This satisfies the "OR" condition of the conclusion.
\end{itemize}
In all cases of the negated premise, at least one of the expressions ($xy$ or $x+y$) evaluates to an odd integer. The contrapositive is proven.
\end{proof}

\newpage
\section*{Untested Concepts}
The following concepts were extensively detailed in the attached reading materials but were \textit{not} explicitly tested in the generated exam variants (to adhere strictly to the structure of the provided previous exam template):
\begin{itemize}
    \item \textbf{Proof by Contradiction (Reductio ad Absurdum):} Utilizing the Axiom of Non-Contradiction to establish impossibility (e.g., proving $\sqrt{2}$ is irrational or the sum of a rational and irrational is irrational).
    \item \textbf{Divisibility Proofs:} Theorems relying heavily on the definition of $a|b \Rightarrow b = ak$ (e.g., proving $a|b \text{ and } a|c \Rightarrow a|(bx+cy)$).
    \item \textbf{Advanced Partition Theory:} Calculating discrete Bell numbers computationally, formal definitions of pairwise disjointness covering a universal set beyond a simple roster example, and Cartesian products of set partitions.
    \item \textbf{Real-World Modeling / Algorithmic Logic:} Applying boolean logic specifically to Quantitative Finance settings (e.g., Binomial Options Pricing constraints, Credit Default Swap triggering logic, structural risk management algorithms).
    \item \textbf{Evaluating Flawed Proofs:} The "Find the Flaw" diagnostic sections involving finding Converse Errors, Inverse Errors, Variable Aliasing, or Division by Zero inside student proofs. 
    \item \textbf{Principle of Inclusion-Exclusion:} Calculating union cardinality algebraically for 3 or more sets using $|A \cup B \cup C| = |A| + |B| + |C| - |A \cap B| - |A \cap C| - |B \cap C| + |A \cap B \cap C|$.
\end{itemize}

\end{document}